\answer{ex:06:1}{$e=(1-e^{-T_a/\tau_e})e^{-d/\tau_e}$。(a)~$0.110$対$4.5\cdot10^{-5}$で、$\approx2.5\cdot10^3$倍。(b)~$\tau_e^*=T_a/\ln(1+T_a/d)=0.59\,\text{s}$。}
\answer{ex:06:2}{速度は毎秒$\nu_x\nu_y(A_+\tau_+-A_-\tau_-)=0.8A_+$、1時間で$\approx2900A_+$。$\tau_-=\tau_+/0.6=33\ms$。}
\answer{ex:06:3}{$\mu^2+s_1^2>s_2^2$のとき$\mathbf e_1$。ここでは$1.01>0.25$で、$\mathbf e_2$方向の散らばりが$25$倍大きいにもかかわらず、ニューロンは$\mathbf e_1$を学ぶ。規則が見つけるのは共分散行列ではなく相関行列の主固有ベクトルである。}
\answer{ex:06:4}{ランプとジュースの間で$V=Re^{-(t_c+L-t)/\tau_\gamma}$なので$\dot V=V/\tau_\gamma$。バーストの面積は$Re^{-L/\tau_\gamma}$で、$0.61R$と$0.78R$。}
\answer{ex:06:5}{SN比は$1/(N+1)$、試行数は$\propto N+1$。$5\cdot10^5$試行$\approx1$か月、$5\cdot10^{11}$試行$\approx8\cdot10^4$年。}
\answer{ex:06:6}{(a)~$k_2=1/\tau_d$、$k_1=1/\tau_d-1/\tau_\gamma$。(b)~偽の誤差は$-(V/\tau_\gamma)\,\varepsilon/(1+\varepsilon)$、$\varepsilon=\tau_d/\tau_\gamma$で、$\tau_d\le0.105\,\text{s}$。}
\answer{ex:06:7}{$0.46$、$0.90$、$0.99$、$0.95$、$d=3\,\text{s}$で$0.70$。点はトレースの割合$F=(1-e^{-T_{\text{cue}}/\tau_e})e^{-d/\tau_e}$の一本の曲線に乗る。$F\gtrsim0.1$で学習し、$F<10^{-2}$でランダムな選択になる。}
\answer{ex:06:8}{$\eta^*=1/\bigl(2\sigma^2(N+2)\bigr)$で、$N+2$試行。$N$個中$m$個がゆらぐ場合は$N(m+2)/m\ge N+2$試行で、スパース性は役に立たない。}
